\chapter{机械波}

振动状态在空间中的传播过程称为\uwave{波动}，简称为\uwave{波}，振动是一个物体的运动，波动是一系列物体的运动，是常见的物质运动形式，我们通常将波动分为以下两大类。
\begin{itemize}
    \item \uwave{机械波}，机械振动在介质中的传播，例如声波，地震波，浅水波，深水波。
    \item \uwave{电磁波}，电磁振荡通过变化的电场和磁场在空间中的传播，例如光波，无线电波。
\end{itemize}
机械波和电磁波虽然在产生和传播机制上有着原理上的本质区别，但是都具有波动的共同特征，都具有一定的传播速度，都伴随着能量的传播，都能产生反射、折射、衍射等现象，同时波动方程具有相同的数学形式。本章以机械波为例，讨论波的特征和基本规律。

\section{机械波的概念}

\subsection{机械波的产生条件}
在弹性介质中，各个质点之间都存在着弹性恢复力的作用当介质中的某一质点发生振动时，就会引起周围其他质点的振动。各个质点间的振动相互联系，由近及远的传播，形成波动。

由此可见，机械波的形成和传播必须具备两个条件
\begin{enumerate}
    \item 能引起机械波的振动的\uwave{波源}。
    \item 能将振动状态传播出去的\uwave{弹性介质}。
\end{enumerate}
在机械波的传播过程中，介质中各质点只在各自的平衡位置附近往复振动，介质中的质点本身并不随机械波传递。即机械波传递质点的运动状态，而不传递质点。\vspace{-10pt}

\subsection{波的类型}
依据质点的振动方向与波的传播方程的关系，可以将波分为两类
\begin{enumerate}
    \item 质点振动方向与波的传播方向相互垂直的波，称为\uwave{横波}。
    \item 质点振动方向与波的传播方向相互平行的波，称为\uwave{纵波}。
\end{enumerate}
例如，声波是纵波，电磁波是横波，地震时产生的地震波则同时具有纵波和横波。

在连续弹性介质中的波动较为抽象，我们以下面的模型直观展示横波和纵波的关系，设想一根弹性绳上等距的穿有若干小球，弹性绳一段固定，而另一端握在手中。

如果手上下振动，此时在绳上可以观察到交替出现的\uwave{波峰}和\uwave{波谷}，这就是横波。

如果手左右推拉，此时在绳上可以观察到交替出新的\uwave{疏部}和\uwave{密部}，这就是纵波。

\begin{Figure}{波的质点分布图}
    \subfigure[横波]
    {
        \label{图:横波}
        \inputfigure[scale=1.2]{Chapter07A_Figure02}
    }
    \subfigure[纵波]
    {
        \label{图:纵波}
        \inputfigure[scale=1.2]{Chapter07A_Figure01}
    }
\end{Figure}
波源做简谐振动时，介质中各点将波源这种简谐振动状态向周围扩散形成的波称为\uwave{简谐波}，在\Refx{图}{横波}和\Refx{图}{纵波}中所示的就是\uwave{简谐横波}和\uwave{简谐纵波}。任何复杂的波，都可以由不同频率不同振幅的简谐波叠加而成，因此简谐波是最简单且又最重要的波。

应当指出，尽管在\Refx{图}{波的质点分布图}中，尽管横波和纵波的形态看起来完全不同，后者具有疏密，前者具有波峰和波谷，但这样的差异只是因为使用了过分具象的质点分布图。

从数学上考察，将横轴记为$x$，将纵轴表示为原先在$x$处的质点相较平衡位置的位移$y$，选取特定时刻$t$观察，那么无论是横波还是纵波，绘制出的图像都是正弦函数，这称为\uwave{波形图}。
\begin{Figure}{波形图}
    \inputfigure[scale=1.2]{Chapter07A_Figure03}
\end{Figure}

对于横波，波形图的纵轴反映了质点在垂直方向的位移，因此与其质点分布图是重叠的。

对于纵波，波形图的纵轴反映了质点在水平方向的位移，因此与其质点分布图差异较大。

在纵波的质点分布图中，波形图为正的点向右偏，波形图为负的点向左偏，因此
\begin{itemize}
    \item 对于密部，其左侧的点向右，其右侧的点向左，这在波形图上表现为“左正右负”，即纵波的密部，对应波形图中，在由波峰到波谷的下降过程中与$x$轴的交点处。
    \item 对于疏部，其左侧的点向左，其右侧的点向右，这在波形图上表现为“左负右正”，即纵波的疏部，对应波形图中，在由波谷到波峰的上升过程中与$x$轴的交点处。
\end{itemize}
因此纵波的“密部疏部”和横波的“波峰波谷”并不是完全对应的。

\subsection{波的传播}
上述讨论中的波动以穿有小球的弹性绳为例，该具象模型可以视为\textbf{波在一维细杆中的传播}，但是一般的波是在空间连续介质中传播的，为了描述这样的波动过程，引入以下的概念。
\begin{itemize}
    \item 我们沿波的传播方向绘制一些带有箭头的线，称为\uwave{波射线}，简称\uwave{波线}。

    \item 我们将振动相位相同的空间各点构成的曲面，称为\uwave{波阵面}，简称\uwave{波面}。
\end{itemize}
波线和波面均可以绘制出任意多个，通常我们在图中只画几个作为代表。

波线在各向同性的介质中，总是垂直通过每一波面。

运用以上概念，一维细杆中的波线即杆所在直线，一维细杆的波面即杆的截面。当波在空间连续各向同性的介质中传播时，可以依据波阵面的形状，将波分为球面波和平面波。
\begin{itemize}
    \item \uwave{球面波}，指的是波阵面为球面的波，其波线是以波源为中心的径向射线。
    
    球面波的波源可以想象为一个体积不断膨胀收缩的球体。

    \item \uwave{平面波}，指的是波阵面为平面的波，其波线是垂直波阵平面的平行射线。
    
    平行波的波源可以想象为一个左右来回颤动的薄板。
\end{itemize}
\begin{Figure}{球面波与平面波}
    \subfigure[平面波]
    {
        \inputfigure{Chapter07A_Figure04}
        \label{图:平面波}
    }
    \qquad
    \subfigure[球面波]
    {
        \inputfigure{Chapter07A_Figure05}
        \label{图:球面波}
    }
\end{Figure}

我们将最远离波源的波阵面称为\uwave{波前}，代表波动传递的前锋位置。

尽管\Refx{图}{球面波与平面波}中为了绘图的便利，将波面以直线和曲线绘制，但是应明确，两者实际代表的是空间中对应位置的平面和球面。平面波和球面波均是在空间中传播的波，平面波中的“平面”，代表的是其波阵面为平面，而不是指波在平面上传播。

不过平面波和球面波确实也有在二维平面介质中传播的简化版本（应注意区分\textbf{平面波}和在\textbf{平面上传播的波}的表述），此时对应的波面退化为平行直线和同心圆环，例如我们在水中投入一颗石子，围绕石子产生了一圈圈的水波涟漪，这就是平面上的球面波的例子。

\section{机械波的属性}

\subsection{波速}
波动是振动状态在空间中的传播，我们将振动状态沿波线的传播速度称为\uwave{波速}和\uwave{相速}。

波速通常用$u$表示，这是为与介质中质点的振动速度$v$加以区分，两者是完全不同的概念，波速$u$衡量了振动状态的传播快慢，振动速度$v$则是质点在其平衡位置附近的速度。

振动状态的传播，是通过弹性介质中质元间的弹性力实现的，因此波速是完全取决于弹性介质的性质，而与波源的属性无关，为了定量计算波速，我们需要引入部分\uwave{弹性力学}的概念。

弹性介质可以视为\uwave{弹性体}，弹性体的研究与先前讨论刚体时，将物体的一切形变忽略的做法是不同的。弹性体在外力作用下，弹性体内部同时产生相应的弹性形变和弹性恢复力。
\begin{itemize}
    \item 这种弹性形变，称为\uwave{应变}，即stress。
    \item 这种弹性恢复力，称为\uwave{应力}，即strain。
\end{itemize}
应力是物体中各部分之间相互作用的内力，但应力实际上并不是力。当我们讨论物体中某处的内力时，就需要设想在该处有一处假象的截面$\delt{S}$。截面$\delt{S}$将弹性物体划为$m_1,m_2$两部分，设$m_1,m_2$间的作用力与反作用力为$\delt{\vb*{f}}$和$-\delt{\vb*{f}}$，定义截面上的应力为以下矢量
\begin{BoxDefinition}[应力]
    \vb*{\tau}=\Lim{\delt{S}}{0}{\frac{\delt{\vb*{f}}}{\delt{S}}}=\Fdif{\vb*{f}}{S}
\end{BoxDefinition}
应力可以认为是物体的单位截面上的内力大小，类似于压强的概念。

应力在国际单位制中的单位是\si{Pa=N\cdot m^{-2}}，称为\uwave{帕斯卡}，量纲是\si{[L^{-1}MT^{-2}]}。

在弹性体的一个截面上的应力$\vb*{\tau}$一般不与此截面垂直，我们可以将其分解为两个分量
\begin{itemize}
    \item 应力在截面上的法向分量$\tau_{\perp}$，通常称为\uwave{正应力}，纯正应力将导致\uwave{体应变}（也称\uwave{容变}）。
    \item 应力在截面上的切向分量$\tau_{\textbf{\scriptsize //}}$\hspace{2.5pt}，通常称为\uwave{剪应力}，纯剪应力将导致\uwave{剪应变}（也称\uwave{剪变}）。
\end{itemize}
对于固体，固体种两种应力都有可能存在，故固体同时可以发生体应变和剪应变。

对于流体，流体内部只存在各向同性的正应力（称为\uwave{压力}或\uwave{张力}），故流体只能发生体应变。

应力的作用效果是应变。那么如何衡量应变的大小？同时应变与应力间存在何种关系？

体应变$\varepsilon_{\text{体}}$表现为弹性体的体积被压缩，以体积的相对变化衡量
\begin{BoxDefinition}[体应变]
    \varepsilon_{\text{体}}=\frac{\delt{V}}{V}
\end{BoxDefinition}
剪应变$\varepsilon_{\text{剪}}$表现为方块在对角挤压下变为菱形，以\uwave{剪变角}衡量
\begin{BoxDefinition}[剪应变]
    \varepsilon_{\text{剪}}=\theta
\end{BoxDefinition}
\begin{Figure}{体应变与剪应变}
    \subfigure[体应变]
    {
        \inputfigure{Chapter07A_Figure06}
    }
    \qquad\qquad
    \subfigure[剪应变]
    {
        \inputfigure{Chapter07A_Figure07}
    }
\end{Figure}
在弹性限度内，正应力$\tau_{\perp}$与体应变成正比，比例系数$K$称为\uwave{体积模量}
\begin{BoxEquation}[容变与体积模量]
    \tau_{\perp}=K\varepsilon_{\text{体}}
\end{BoxEquation}
在弹性限度内，剪应力$\tau_{\textbf{\scriptsize //}}$\hspace{2.5pt}与剪应变成正比，比例系数$G$称为\uwave{剪变模量}
\begin{BoxEquation}[剪变与剪变模量]
    \tau_{\textbf{\scriptsize //}}\hspace{2.5pt}=G\varepsilon_{\text{剪}}
\end{BoxEquation}
\begin{Figure}{杆的拉伸压缩形变}
    \inputfigure{Chapter07A_Figure08}
\end{Figure}
我们再来考量另外一种情况，当长度为$l$的直杆两端加上与杆平行的力$\vb*{f}$的拉伸或压缩时，杆的长度将会有所改变，杆的应变可以用长度的相对变化$\varepsilon=\dfrac{\delt{l}}{l}$衡量，称为\uwave{长变}。

细杆的应力则为$\tau=\dfrac{f}{\delt{S}}$，其中$\delt{S}$为杆的截面面积。

细杆在弹性限度内，应力$\tau$与应变$\varepsilon$成正比，比例系数$E$称为\uwave{杨氏模量}
\begin{BoxEquation}[长变与杨氏模量]
    \tau=E\varepsilon=E\frac{\delt{l}}{l}
\end{BoxEquation}
胡克首先研究并发表了弹性细杆拉伸压缩形变的规律\Refx{公式}{长变与杨氏模量}，现在将此式以及在胡克之后的时代种发现的包括\Refx{公式}{容变与体积模量}以及\Refx{公式}{剪变与剪变模量}在内的“应变与应力成正比的规律”统称为\uwave{胡克定律}。

细杆在\uwave{纵向形变}的同时，也总伴有\uwave{横向形变}，将横向形变和纵向形变的比称为\uwave{泊松比}。

细杆在单向拉伸压缩的长变过程，事实上可以证明是容变和剪变的结合，因而我们引入的三个弹性模量$K,G,E$和泊松比之间是存在联系的的，根据弹性理论可以证明这四个物理量之中只有两个是独立的。我们不关心具体的数学关系，关键是理解长变是容变和剪变的结合。

机械纵波的传递需要正应力的作用，机械横波的传递需要剪应力的作用。

流体中只存在正应力，因此在流体中只能传递纵波，由于正应力的作用结果是容变，而容变又与体积模量$K$有关，因此流体中的波速可以由流体的体积模量$K$计算
\begin{BoxEquation}[机械纵波在流体中的波速]
    u=\sqrt{\frac{K}{\rho}}
\end{BoxEquation}
固体中同时存在正应力和剪应力，因此在固体中可以同时传递纵波和横波，但是固体的体积是很难改变的，因此单纯的容变较难发生，实际上当固体传递纵波时，其发生的是长变。

固体传递纵波时，发生长变，纵波波速与杨氏模量有关
\begin{BoxEquation}[机械横波在固体中的波速]
    u=\sqrt{\frac{E}{\rho}}
\end{BoxEquation}
固体传递横波时，发生剪变，横波波速与剪变模量有关
\begin{BoxEquation}[机械纵波在固体中的波速]
    u=\sqrt{\frac{G}{\rho}}
\end{BoxEquation}
以上的$\rho$均表示流体介质和固体介质的密度。

弹性模量$E,K,G$的量纲均与应力相同，这是因为关注到应变$\varepsilon$总是无量纲值。

弹性模量$E,K,G$无论固体软硬，均具有大致相同的数量级，大约在$10^{10}$\si{Pa}左右的数量级，即相当于$10^5$倍的大气压，这一结果的原理需要通过量子理论来解释，此处从略。

在地震过程产生的地震波中，纵波（P波）首先到达地表，横波（S波）随后才到达，由此可以根据这一波速关系大致推断出地壳的杨氏模量大于剪变模量，即$E>G$。

在地震中，纵波将导致建筑物上下晃动，横波将导致建筑物左右晃动，由于现代的高层建筑垂直耸立，普遍难以耐受横波的水平拉扯，同时横波的振幅往往也更大，因此地震中造成较为严重破坏的往往是横波，现代地震预警利用的就是纵波和横波间数十秒的时间差。

\subsection{波长}
波动是振动状态在空间中的传播，而振动相位的传播是受到波速的限制。

设波速为$u$，设波源的振动周期是$T$，在$t=0$时波源由平衡位置起振，当$t=T$时波源处再次回到初始相位，而此时$t=0$时的相位传播到$x=uT$处，由$x=uT$至波源$x=0$的质点分别具有$0$至$2\pi$的相位，即两者之间具有一个完整的波形。

我们将波动一个完整波形的产生所需的时间称为\uwave{波的周期}，仍用$T$表示。

我们将波动一个完整波形的长度称为\uwave{波长}，使用$\lambda$表示。

根据定义，波的周期$T$与振动周期完全相同，因此类似的也有波的频率$\nu$和波的角频率$\omega$的概念，两者均与振动的频率和角频率相等，也都满足$T=\dfrac{1}{\nu}$以及$T=\dfrac{2\pi}{\omega}$的公式。

根据以上的讨论，波长和波的周期间存在以下关系
\begin{BoxEquation}[波长和周期的关系]
    \lambda=uT
\end{BoxEquation}
由此可见，波长是一个周期内波的传播距离，波的周期是波传播一个波长所需的时间。

类似于\Refx{定义}{频率}，定义波长的倒数为\uwave{波数}，使用$\widetilde{\nu}$表示
\begin{BoxDefinition}
    \widetilde{\nu}=\frac{1}{\lambda}
\end{BoxDefinition}
类似于\Refx{定义}{角频率}，定义波长的倒数的$2\pi$倍为\uwave{角波数}，使用$k$表示
\begin{BoxDefinition}
    k=\frac{2\pi}{\lambda}
\end{BoxDefinition}
由于周期$T$描述了传递一个波形的时间，同时波长$\lambda$描述了一个波形的空间长度
\begin{itemize}
    \item 角频率描述了单位时间内的相位变化，频率描述了单位时间内的振动次数。
    \item 角波数描述了单位空间上的相位变化，波数描述了单位空间上的波形个数。
\end{itemize}
空间周期是时间周期的波速倍$\lambda=uT$，时间频率是空间频率的波速倍，即$\nu=u\widetilde{\nu},~\omega=uk$。

时间概念组“周期$T$~-~频率$\nu$~-~角频率$\omega$”表示了波的\uwave{时间周期}和\uwave{时间频率}。

空间概念组“波长$\lambda$\hspace{2pt}~-~波数$\widetilde{\nu}$~-~角波数$k$”表示了波的\uwave{空间周期}和\uwave{空间频率}。

波速$u$是沟通时间概念组和空间概念组的桥梁。

角频率和角波数分别代表单位时间和单位空间的相位变化，这就是说，时间间隔和空间距离都可以影响振动相位，因此，时间和空间在波动中可以通过相位的度量相互等效。

角波数和波数在国际单位制中的单位是\si{m^{-1}}，量纲是\si{[M^{-1}]}。

\section{平面简谐波的表达式}
描述波动传播规律的数学函数式，称为\uwave{波函数}或\uwave{波动表达式}，由于波动是介质中一系列物质的运动，因此描述波动的波函数，是关于时间和空间的二元函数。

我们以最简单的波源振动\textbf{简谐振动}和最简单的波阵面\textbf{平面波}，构成的\uwave{平面简谐行波}为例，推导波函数。其中的\uwave{行波}是对在空间中运动的波，与稍后的\uwave{驻波}相对应。

\subsection{波函数的推导}
设一个无限大，均匀，无吸收，各向同性的理想介质，设一个平面在该介质中振动，沿垂直平面的方向产生了一个以速度$u$沿$x$轴正向传播的的平面波，由于平面波的各波线的波动状况是完全相同，我们任选其中一根波线进行研究，并设振源处的坐标为$x=0$。

波函数是一个关于空间$x$和时间$t$，波函数的意义是位于$x$处的质点在$t$时的振动状况，即相对平衡位置的位矢，先前我们在振动中总是使用$x$表述质点相对平衡位置的位矢，但是考虑到波动中$x$已被用于表示质点的平衡位置，因此改用$y$表示相对平衡位置的位矢。

应当指出，这里的$x,y$记号并未隐式假定两者的方向关系
\begin{itemize}
    \item 对于横波，振动方向与波的传播方向垂直，那么振动$y$的方向就与位置$x$垂直。
    \item 对于纵波，振动方向与波的传播方向平行，那么振动$y$的方向就与位置$x$平行。
\end{itemize}
由此可见，波函数无关横纵波，波函数仅关心偏离平衡位置的位矢的大小，不考虑方向。

设波源$O$的振动方程为
\begin{Equation}
    y=A\cos(\omega t+\phi)
\end{Equation}
现在考察沿波线距离波源为$x$的任意一点$P$的振动方程，由于波动是振动状态的空间传播，因此$P$点在$t$时的相位，实际上是由$O$点在$t-\dfrac{x}{u}$历经$\delt{t}=\dfrac{x}{u}$的时间传播而来的。

这就得到了空间中任意一点的振动方程，这便是\uwave{平面简谐行波的波函数}
\begin{BoxEquation}[波函数]
    y=A\cos\qty[\omega\qty(t-\frac{x}{u})+\phi]
\end{BoxEquation}

\subsection{波函数的物理意义}
\Refx{公式}{波函数}是通过将空间距离等效为时间差距得到的，但是我们也可以将其化为对称的形式，展开并代入角频率和角波数的关系$\omega=uk$，得到
\begin{equation*}
    y=A\cos\qty[\omega\qty(t-\frac{x}{u})+\phi]=A\cos\qty[\omega t-\frac{\omega}{u}x+\phi]=A\cos\qty[\omega t-kx+\phi]
\end{equation*}\vspace{-15pt}
\begin{BoxEquation}[波函数的频率表示]
    y=A\cos\qty[\omega t-kx+\phi]
\end{BoxEquation}
\Refx{公式}{波函数的频率表示}指出，空间和时间可以以对等的方式影响振动相位
\begin{itemize}
    \item 时间对相位具有正贡献，时间与角频率的积$\omega t$的\uwave{正值}是时间造成的相位差，也就是说，较晚的时间意味着更为领先的相位，时间差$\delt{t}$的领先意味着$\omega\delt{t}$的相位领先。
    \begin{center}
        沿时间演化方向，每个质点的振动状态逐渐领先。
    \end{center}
    \item 空间对相位具有负贡献，空间与角波数的积$kt$的\uwave{负值}是时间造成的相位差，也就是说，较远的空间意味着更为落后的相位，空间差$\delt{x}$的领先意味着$kt$的相位差落后。
    \begin{center}
        沿空间传播方向，各个质点的振动状态依次落后。
    \end{center}
\end{itemize}
以上讨论指出，时间对相位具有正贡献，空间对相位具有负贡献，但这仅仅是对于该形式的函数而言，因为相位的正负本身具有相对性，如果利用$\cos(\theta)=\cos(-\theta)$反转相位的正负，那么结论将刚好相反，然而不变的是，时间和空间对相位的作用总是相反的。
\begin{Figure}{波形图与振动图}
    \subfigure[波形图]
    {
        \inputfigure[scale=1]{Chapter07B_Figure02}
    }\quad\quad
    \subfigure[振动图]
    {
        \inputfigure[scale=1]{Chapter07B_Figure01}
    }
    \vspace{-15pt}
\end{Figure}\vspace{-10pt}
当$t$\hspace{4pt}取定时，关于空间$x$的函数$y=A\cos[kx-(\phi+\omega t)]$的$y-x$图称为$t$\hspace{4pt}时的\uwave{波形图}。

当$x$取定时，关于时间$t$\hspace{4pt}的函数$y=A\cos[\omega t+(\phi-kx)]$的$y-t$\hspace{5pt}图称为$x$处的\uwave{振动图}。
\begin{itemize}
    \item 波形图关于空间，以空间方向为相位正方向，在这种情况下时间对相位的影响是负的，因此时间上领先的\hspace{4.5pt}$t+\delt{t}$\hspace{4.5pt}时的波形图在相位上落后于\hspace{4.5pt}$t$\hspace{4.5pt}时的波形图$\omega\delt{t}$。
    \item 振动图关于时间，以时间方向为相位正方向，在这种情况下空间对相位的影响是负的，因此空间上领先的$x+\delt{x}$处的振动图在相位上落后于$x$处的波形图$k\delt{x}$。
\end{itemize}
在\Refx{图}{波形图}中，波速的大小表现为波形曲线随时间的运动快慢，波速的方向即波形运动的方向，该例子中当时间由$t$增加至$t+\delt{t}$，波形曲线向右（由蓝至红），波速故向右。

\Refx{公式}{波函数的频率表示}中，代入周期$T=\dfrac{2\pi}{\omega}$和波长$\lambda=\dfrac{2\pi}{k}$
\begin{BoxEquation}[波函数的周期表示]
    y=A\cos\qty[2\pi\qty(\frac{t}{T}-\frac{x}{\lambda})+\phi]
\end{BoxEquation}
\Refx{公式}{波函数的周期表示}指出，波动同时具有空间周期性和时间周期性
\begin{itemize}
    \item 当时间$t$相差$T,2T,3T,\cdots,nT$时，同一位置的质点的状态相同，即\uwave{波的时间周期性}。
    \item 当空间$x$相隔$\lambda,2\lambda,3\lambda,\cdots,n\lambda$时，同一时间的质点的状态相同，即\uwave{波的空间周期性}。
\end{itemize}
对\Refx{公式}{波函数的频率表示}求时间的一阶偏导，即质点的振动速度
\begin{BoxEquation}[波的振动速度]
    v=\Fpif{y}{t}=-A\omega\sin\qty[\omega t-kx+\phi]   
\end{BoxEquation}
对\Refx{公式}{波函数的频率表示}求时间的二阶偏导，即质点的振动加速度
\begin{BoxEquation}[波的振动加速度]
    a=\Fpif{y}{t}^2=-A\omega^2\sin\qty[\omega t-kx+\phi]
\end{BoxEquation}
注意区分波速$u$与波在传播中质点的振动速度$v$的差异。

\subsection{平面波的波动方程}
求\Refx{公式}{波函数的频率表示}对时间的二阶偏导数
\begin{Equation}[波对时间二阶导]
    \Fpif{y}{t}^2=-A\omega^2\sin\qty[\omega t-kx+\phi]
\end{Equation}
求\Refx{公式}{波函数的频率表示}对空间的二阶偏导数
\begin{Equation}[波对空间二阶导]
    \Fpif{y}{x}^2=-Ak^2\sin\qty[\omega t-kx+\phi]
\end{Equation}
比较\ref{式:波对时间二阶导}与式\ref{式:波对空间二阶导}（将两者做比），考虑到$\omega^2=u^2k^2$
\begin{BoxEquation}[平面波的波动方程]
    \Fpif{y}{x}^2=\frac{1}{u^2}\Fpif{y}{t}^2
\end{BoxEquation}
上式称为\uwave{平面波的波动方程}，这是一个偏微分方程，尽管是由平面简谐波的波函数推导而来，但事实上该偏微分方程的解远不止简谐波，其反映了各向同性的均匀介质中的一切平面波的共同特征“振动位矢对空间的二阶偏导，等同于振动位矢对时间的二阶偏导除以波速的平方”，其解是平面波的波函数的全体，而平面简谐波的波函数是其中的一个解。

在这一章中，我们对波动的研究均是基于这样一个假设，即振动形式在弹性介质中会以波的方式传播，尽管这个假设似乎十分符合直观，但是客观规律并不以人的意志为转移，因此我们需要通过动力学的观点，论证弹性介质的动力学规律必然会导致波动的运动学结果。

因此我们的目的就是要推导出\Refx{公式}{平面波的波动方程}。

下面以固体细长棒中的平面纵波为例，设棒的截面为$S$，设棒的密度为$\rho$。
\begin{Figure}{纵波在细棒中的传播}
    \inputfigure[scale=0.85]{Chapter07B_Figure03}
\end{Figure}
当纵波传播时棒会发生形变，棒中不同部位被不同程度的拉伸和压缩，以棒长方向为$x$轴，在棒中任取长为$\dx$的质元，设微元的左截面和右截面分布位于$x$和$x+\dx$处。

我们知道，对于函数$f(x)$，有$f(x+\dx)=f(x)+\dif{f}(x)+o(\dx)$，而$o(\dx)$在物理上使用微元法时是可以忽略不计的，但在这里力和振动位移是关于$x,t$的二元函数$F(x,t),y(x,t)$，因此当保持$t$不变，考察$(x+\dx,t)$处的函数值时，要改用对$x$的偏微分表示，即
\begin{equation*}
    F(x+\dx,t)=F+\Fpif{F}{x}\dx\qquad\quad y(x+\dx,t)=y+\Fpif{y}{x}\dx
\end{equation*}
关注到质元$\dx$受到的合力为$\Fpif{F}{x}\dx$，方向向右，根据\Refx{定律}{牛顿第二定律}
\begin{Equation}[波动方程1]
    \Fpif{F}{x}\dx=a\dif{m}=a\rho S\dx
\end{Equation}
代入\Refx{定义}{瞬时加速度}
\begin{Equation}[波动方程2]
    \Fpif{F}{x}\dx=\Fpif{y}{t}^2\rho S\dx
\end{Equation}
根据\Refx{定义}{应力}，由于此处的细杆是均匀的，因此有$F=\tau S$，故
\begin{Equation}[波动方程3]
    \Fpif{F}{x}=\Fpif{\tau}{x}S
\end{Equation}
在式\ref{式:波动方程2}中代入式\ref{式:波动方程3}，并约去$S\dx$
\begin{Equation}[波动方程4]
    \Fpif{\tau}{x}=\rho\Fpif{y}{t}^2
\end{Equation}
根据\Refx{公式}{长变与杨氏模量}，此处有$\tau=E\varepsilon=E\dfrac{\delt{l}}{l}$，其中质元的原长$l=\dx$，而质元的长度增量$\delt{l}=\Fpif{y}{x}\dx$，从而将两者相除将$\dx$约去，这就得到$\tau=E\Fpif{y}{x}$。

从而就可以求出应力$\tau$对空间$x$的偏导
\begin{Equation}[波动方程5]
    \Fpif{\tau}{x}=\Fpif{x}\qty(E\Fpif{y}{x})=E\Fpif{y}{x}^2
\end{Equation}
将式\ref{式:波动方程5}代入式\ref{式:波动方程4}，这就得到
\begin{Equation}[波动方程6]
    \Fpif{y}{x}^2=\frac{\rho}{E}\Fpif{y}{t}^2
\end{Equation}
如果我们令$u=\sqrt{\dfrac{E}{\rho}}$，那么式\ref{式:波动方程6}中就有$\dfrac{\rho}{E}=\dfrac{1}{u^2}$，这就得到了\Refx{公式}{平面波的波动方程}中的偏微分方程，这就证明了弹性介质的动力学性质必然导致波动的运动学结果。

此处的$u$仅仅是一个代换变量，但是通过波动方程解的物理意义可知，这里的$u$实际上就是我们所熟知的波速，这一代换公式也是\Refx{公式}{机械纵波在固体中的波速}的来源。

此处以纵波为例推导，但事实上\Refx{公式}{平面波的波动方程}普遍适用于纵波和横波。